\section{Method}\label{method}

\subsection{A world compiler}\label{a-world-compiler}

SCORE compiles a world once, offline, in batches; nothing is generated while it is played, and each scale of a world (a room, a planet's ground, an orbit) is its own kind of world. Table 1 lists the stages. Each stage's output is explicit, cached and checked before the next paid stage starts.

\textbf{Table 1.} Stages, their outputs and the checks that gate the next stage.

{\def\LTcaptype{none} % do not increment counter
\begin{longtable}[]{@{}
  >{\raggedright\arraybackslash}p{(\linewidth - 4\tabcolsep) * \real{0.3333}}
  >{\raggedright\arraybackslash}p{(\linewidth - 4\tabcolsep) * \real{0.3333}}
  >{\raggedright\arraybackslash}p{(\linewidth - 4\tabcolsep) * \real{0.3333}}@{}}
\toprule\noalign{}
\begin{minipage}[b]{\linewidth}\raggedright
Stage
\end{minipage} & \begin{minipage}[b]{\linewidth}\raggedright
Output
\end{minipage} & \begin{minipage}[b]{\linewidth}\raggedright
Check before the next stage
\end{minipage} \\
\midrule\noalign{}
\endhead
\bottomrule\noalign{}
\endlastfoot
Concept & picked concept picture & creator's pick \\
Plan & dimensioned plan in metres & room and door checks \\
Inventory & every object as a row; children on parents & concept-density check \\
Close-ups & one clean picture per object & one object per picture \\
Shape & code builder or picture-to-3D model per kind & parts check, model check \\
Surfaces & baked maps from the shared library & palette sweep, storage budget \\
Sound & sound per surface, room and object & loudness and fault checks \\
Assembly & the OpenUSD scene and engine builds & scene check, made-only check \\
Review & shots from fixed cameras & creator's keep or send-back \\
\end{longtable}
}

\textbf{The canonical scene.} All stages write one OpenUSD stage \citep{openusd} with glTF geometry \citep{gltf2}. Every object carries its name, kind, place in metres, collision, the library surface of each part, its sound, and tags for what a person can do with it (door, seat, terminal, airlock). The compiler owns a generated base layer and may rewrite it; the creator's changes live in edit layers above it, so a regeneration replaces the base and keeps the edits. Engines and simulators (Godot, Blender, Unreal, robotics simulators) load the same stage through thin adapters.

\subsection{Checks before spending}\label{checks-before-spending}

A fault costs nothing in the plan, cents at the picture, about €0.18 and 75 minutes at the 3D model, and an evening of the creator's time once a room is built and played. Each paid stage is therefore gated by deterministic checks on real geometry, never by an agent's judgement, and a failure sends work back one stage. Room checks cast rays for leaks through walls and roof, pieces blocking openings and pieces off their surface, in about two seconds a room. A door check requires every door to separate its two sides. A concept-density check requires every element of the concept to have an inventory row or a written reason. A model check requires closed solids with walls of at least 3 mm and fails generated boxes that lean more than 2 degrees. A palette sweep and a storage budget per room follow, and after assembly a scene check finds what floats or sinks and fails on any visible mesh the route did not make.

\subsection{Code or pipeline, one model per object}\label{code-or-pipeline-one-model-per-object}

Picture-to-3D models build chunky solids well and flat panels, thin beams and openings badly; agent-written code builders are exact, light and straight, but show only the parts the agent wrote. SCORE routes each object kind by the \textbf{parts check}: a kind is built in code only when its code build shows every part its clean close-up shows, each at least 10 percent visible and none under a label. Every other kind goes through the prop pipeline: Pixal3D \citep{li2026pixal3d} in a cloud batch, closed into a solid, triangles set by size. Composites are split before generation, one model per real-world object: a tool board is a board carrying its tools, a console its monitors and keyboards, each made on its own.

\subsection{Surfaces and sound}\label{surfaces-and-sound}

Models and code give shape only. Every part takes one surface from a shared library of ProcFunc functions \citep{raistrick2026procfunc} on Infinigen's shaders \citep{raistrick2023infinigen}, coloured from the place's palette, with wear, dirt and seed as named settings. A room has one wear setting, applied from causes: edges, foot traffic, drips. Parts of a generated model come from PartCrafter \citep{lin2025partcrafter}, and each part takes one surface, chosen from its picture. Labels are a few decals on clear flat spots. Sound is compiled the same way: each surface carries its footsteps and impacts, each room its echo from its size and surfaces, each object its own sound, generated in one batch by MOSS-SoundEffect \citep{moss2026}, chosen by CLAP prompt match \citep{wu2022clap} less measured faults, and levelled under a -3 dBTP true peak. The libraries of surfaces and sounds grow with each world and are reused in the next.

\subsection{Cloud batches and human picks}\label{cloud-batches-and-human-picks}

Every heavy step runs in a rented cloud machine that is deleted afterwards, records its cost per batch, and stops itself when its output grows past its expected size. The creator picks at fixed points only. Where two routes compete the pick is blind: both drawn the same way, labelled X and Y at random.

\subsection{The optional world step}\label{the-optional-world-step}

A concept shows one view. An optional world step turns it into a walkable whole-room reference, so close-ups can show every object from every side and the walls the concept does not show can be filled in its style. World Labs Marble served this step in our runs; its output is never shipped, and anything derived from it is credited ``Generated using World Labs''.
